% ============================================================
% 试卷模板 - 通用版
% 使用方法: 复制此文件，根据需要修改配置和内容
% 编译命令: xelatex exam.tex
% ============================================================

\documentclass[11pt,a4paper]{article}

% ============ 页面设置 ============
\usepackage[top=2cm,bottom=2cm,left=2.2cm,right=2.2cm]{geometry}

% ============ 中文支持 ============
\usepackage{ctex}

% ============ 代码高亮（可选，编程类试卷使用）============
\usepackage{listings}
\usepackage{xcolor}

% ============ 列表和枚举 ============
\usepackage{enumitem}

% ============ 页眉页脚 ============
\usepackage{fancyhdr}
\usepackage{lastpage}

% ============ 表格 ============
\usepackage{booktabs}
\usepackage{tabularx}

% ============ 其他 ============
\usepackage{ulem}      % 下划线
\usepackage{tikz}      % 绘图
\usepackage{amssymb}   % 数学符号
\usepackage{pifont}    % 特殊符号（✓ ✗）

% ============ 字体设置 ============
% macOS 字体配置
\setCJKmainfont{Songti SC}[BoldFont={Heiti SC}]
\setCJKsansfont{Heiti SC}
\setmainfont{Times New Roman}[BoldFont={Times New Roman Bold}]
\setmonofont{Menlo}[Scale=0.88]

% Windows 字体配置（取消上面的注释，注释掉 macOS 配置）
% \setCJKmainfont{SimSun}[BoldFont={SimHei}]
% \setCJKsansfont{SimHei}
% \setmainfont{Times New Roman}
% \setmonofont{Consolas}[Scale=0.88]

% Linux 字体配置（取消上面的注释，注释掉其他配置）
% \setCJKmainfont{Noto Serif CJK SC}[BoldFont={Noto Sans CJK SC}]
% \setCJKsansfont{Noto Sans CJK SC}
% \setmainfont{DejaVu Serif}
% \setmonofont{DejaVu Sans Mono}[Scale=0.88]

% ============ 代码样式（可选）============
\definecolor{codebg}{HTML}{F5F5F5}
\definecolor{codeframe}{HTML}{DDDDDD}
\definecolor{keyword}{HTML}{0000FF}
\definecolor{comment}{HTML}{008000}
\definecolor{stringcol}{HTML}{A31515}

% 定义代码语言（以 Python 为例，可根据需要修改）
\lstdefinelanguage{Python}{
  keywords={and,as,assert,break,class,continue,def,del,elif,else,except,
    finally,for,from,global,if,import,in,is,lambda,not,or,pass,raise,
    return,try,while,with,yield,True,False,None},
  sensitive=true,
  morecomment=[l]{\#},
  morecomment=[s]{'''}{'''},
  morecomment=[s]{"""}{"""},
  morestring=[b]",
  morestring=[b]',
}

\lstset{
  language=Python,
  basicstyle=\ttfamily\small,
  keywordstyle=\color{keyword}\bfseries,
  commentstyle=\color{comment}\itshape,
  stringstyle=\color{stringcol},
  backgroundcolor=\color{codebg},
  frame=single,
  rulecolor=\color{codeframe},
  xleftmargin=1em,
  xrightmargin=0.5em,
  framexleftmargin=0.5em,
  numbers=none,
  breaklines=true,
  breakatwhitespace=false,
  tabsize=4,
  showstringspaces=false,
  aboveskip=0.8em,
  belowskip=0.6em,
  escapeinside={(*}{*)},
}

% ============ 页眉页脚设置 ============
\pagestyle{fancy}
\fancyhf{}
\renewcommand{\headrulewidth}{0.4pt}
\renewcommand{\footrulewidth}{0pt}
\fancyfoot[C]{\small 第 \thepage\ 页 / 共 \pageref{exam:last} 页}

% 修改这里的页眉内容
\fancyhead[L]{\small XXX 公司 YYY 岗位笔试题（A 卷）}
\fancyhead[R]{\small\color{red}内部资料 · 禁止外传}

% ============ 列表紧凑 ============
\setlist{nosep,leftmargin=2em}

% ============ 选项列表环境 ============
\newlist{choices}{enumerate}{1}
\setlist[choices]{label=\Alph*.,leftmargin=2.5em,labelsep=0.5em,nosep,topsep=0.3em,itemsep=0.1em}

% ============ 判断题空格 ============
\newcommand{\tfblank}{\makebox[2em]{\rule[-0.3ex]{1.5em}{0.4pt}}}

% ============ 简答题答题区 ============
\newcommand{\answerlines}[1]{%
  \par\vspace{\dimexpr #1\baselineskip + 0.5em\relax}%
}

% ============================================================
% 文档内容开始
% ============================================================
\begin{document}

% ============ 试卷标题 ============
\begin{center}
{\LARGE\bfseries\sffamily XXX 公司 2026 年}
\hspace{0.5em}
{\LARGE\bfseries\sffamily YYY 岗位笔试题（A 卷）}

\vspace{0.6em}
{\large 总分：100 分 \quad 考试时间：100 分钟}

\vspace{0.4em}
\begin{tabularx}{0.7\textwidth}{Xr}
姓名：\underline{\hspace{6em}} & 日期：\underline{\hspace{6em}}
\end{tabularx}
\end{center}

\vspace{0.3em}
\noindent\rule{\linewidth}{0.4pt}
\vspace{0.2em}

% ============ 一、单选题 ============
\subsection*{一、单选题（每题 3 分，共 30 分）}
\vspace{0.2em}
\noindent\small{请将正确选项字母填入题号前的括号内。}
\vspace{0.3em}

\noindent(\hspace{1.5em})\enspace\textbf{1.} 这是第一道单选题的题目？

\begin{choices}
\item 选项 A
\item 选项 B
\item 选项 C
\item 选项 D
\end{choices}

\noindent(\hspace{1.5em})\enspace\textbf{2.} 这是第二道单选题的题目？

\begin{choices}
\item 选项 A
\item 选项 B
\item 选项 C
\item 选项 D
\end{choices}

% ... 继续添加更多单选题 ...

% ============ 二、多选题 ============
\subsection*{二、多选题（每题 4 分，共 20 分，少选错选多选不得分）}
\vspace{0.2em}
\noindent\small{请将所有正确选项字母填入题号前的括号内。}
\vspace{0.3em}

\noindent(\hspace{1.5em})\enspace\textbf{11.} 以下哪些说法是正确的？

\begin{choices}
\item 选项 A
\item 选项 B
\item 选项 C
\item 选项 D
\end{choices}

% ... 继续添加更多多选题 ...

% ============ 三、判断题 ============
\subsection*{三、判断说法正误（每题 2 分，共 20 分，正确打 $\checkmark$，错误打 \ding{55}）}
\vspace{0.2em}
\noindent\small{请在题号前的括号内打 $\checkmark$ 或 \ding{55}。}
\vspace{0.3em}

\noindent(\hspace{1.5em})\enspace\textbf{16.} 这是一个判断题的陈述。

\noindent(\hspace{1.5em})\enspace\textbf{17.} 这是另一个判断题的陈述。

% ... 继续添加更多判断题 ...

% ============ 四、简答题 ============
\subsection*{四、简答题（共 30 分）}
\vspace{0.3em}

\noindent\textbf{26.}（10 分）请简述 XXX 的基本原理。

\par\noindent（1）小问 1（3 分）
\answerlines{4}

\par\noindent（2）小问 2（4 分）
\answerlines{5}

\par\noindent（3）小问 3（3 分）
\answerlines{4}

\noindent\textbf{27.}（10 分）阅读以下代码，回答问题。

% 如果有代码题，使用 lstlisting 环境
\begin{lstlisting}
def example():
    print("Hello, World!")
\end{lstlisting}

\par\noindent（1）这段代码的输出是什么？（3 分）
\answerlines{3}

\par\noindent（2）请解释代码的工作原理。（7 分）
\answerlines{6}

\noindent\textbf{28.}（10 分）请比较 A 和 B 两种方法的优缺点。

\answerlines{10}

% ============ 试卷结束标记 ============
\label{exam:last}

% ============================================================
% 参考答案（可选，可放在试卷末尾或单独文件）
% ============================================================
\newpage
\setcounter{page}{1}
\fancyfoot[C]{\small 参考答案 第 \thepage\ 页 / 共 \pageref{ans:last} 页}

\subsection*{参考答案（A卷）}

\subsubsection*{一、单选题}
1. \textbf{A} — 解析说明。

2. \textbf{C} — 解析说明。

% ... 继续添加答案 ...

\subsubsection*{二、多选题}
11. \textbf{A, B, D} — 解析说明。

% ... 继续添加答案 ...

\subsubsection*{三、判断题}
16. \textbf{$\checkmark$} — 解析说明。

17. \textbf{\ding{55}} — 解析说明。

% ... 继续添加答案 ...

\subsubsection*{四、简答题}

\textbf{26.}

（1）答案要点：
\par\noindent\hspace{1em}$\bullet$ 要点 1
\par\noindent\hspace{1em}$\bullet$ 要点 2

（2）答案要点：
\par\noindent\hspace{1em}$\bullet$ 要点 1
\par\noindent\hspace{1em}$\bullet$ 要点 2

（3）答案要点：
\par\noindent\hspace{1em}$\bullet$ 要点 1

\textbf{27.}

（1）输出结果。

（2）工作原理解释。

\textbf{28.}

比较分析要点。

\label{ans:last}

\end{document}
